% --- Chapter: Control Flow ---
\section{Control Flow}

\subsection{Conditional Statements}

CRISP supports \texttt{if}, \texttt{else if}, and \texttt{else} with block bodies:

\begin{tcolorbox}[colback=codebg, colframe=darkpurple, colbacktitle=darkpurple, title=\textbf{\textcolor{codegold}{Conditional Statements}}, fonttitle=\bfseries, sharp corners]
\begin{lstlisting}
# Basic if/else
let x = 42;
if x > 100 {
    say "Large";
} else {
    say "Small";
}

# if/else if/else chain
let score = 85;
if score >= 90 {
    say "A";
} else if score >= 80 {
    say "B";
} else if score >= 70 {
    say "C";
} else {
    say "F";
}

# Single-expression conditions
let name = "CRISP";
if name == "CRISP" {
    say "It's CRISP!";
}
\end{lstlisting}
\end{tcolorbox}

\subsection{While Loops}

The \texttt{while} loop evaluates a condition before each iteration:

\begin{tcolorbox}[colback=codebg, colframe=darkpurple, colbacktitle=darkpurple, title=\textbf{\textcolor{codegold}{While Loops}}, fonttitle=\bfseries, sharp corners]
\begin{lstlisting}
# Basic while loop
let i = 0;
while i < 5 {
    say i;
    i = i + 1;
}
# Output: 0 1 2 3 4

# Infinite loop with break
let count = 0;
while true {
    say count;
    count = count + 1;
    if count >= 5 { break; }
}
# Output: 0 1 2 3 4
\end{lstlisting}
\end{tcolorbox}

\subsection{For Loops}

The \texttt{for} loop iterates over arrays, hashes, strings, and ranges:

\begin{tcolorbox}[colback=codebg, colframe=darkpurple, colbacktitle=darkpurple, title=\textbf{\textcolor{codegold}{For Loops}}, fonttitle=\bfseries, sharp corners]
\begin{lstlisting}
# For loop over array
let @fruits = ["apple", "banana", "cherry"];
for fruit in @fruits {
    say fruit;
}
# Output: apple banana cherry

# For loop over range
for i in 0..5 {
    say i;
}
# Output: 0 1 2 3 4

# For loop over hash (iterates values)
let %ages = { alice => 25, bob => 30 };
for age in %ages {
    say age;
}
# Output: 25 30

# For loop over string (iterates characters)
for ch in "CRISP" {
    say ch;
}
# Output: C R I S P
\end{lstlisting}
\end{tcolorbox}

\subsection{Break and Continue}

\texttt{break} exits the loop immediately. \texttt{continue} skips to the next iteration:

\begin{tcolorbox}[colback=codebg, colframe=darkpurple, colbacktitle=darkpurple, title=\textbf{\textcolor{codegold}{Break and Continue}}, fonttitle=\bfseries, sharp corners]
\begin{lstlisting}
# Break — exit loop early
let i = 0;
while true {
    say i;
    i = i + 1;
    if i >= 5 { break; }
}
# Output: 0 1 2 3 4

# Continue — skip to next iteration
for i in 0..10 {
    if i % 2 == 0 { continue; }  # skip evens
    say i;
}
# Output: 1 3 5 7 9

# Both together
for i in 0..10 {
    if i == 2 { continue; }
    if i == 7 { break; }
    say i;
}
# Output: 0 1 3 4 5 6
\end{lstlisting}
\end{tcolorbox}

\subsection{Control Flow and Try/Finally}

\texttt{break} and \texttt{continue} propagate correctly through \texttt{try/finally} — the \texttt{finally} block always executes before the signal propagates:

\begin{tcolorbox}[colback=codebg, colframe=darkpurple, colbacktitle=darkpurple, title=\textbf{\textcolor{codegold}{Break with Finally}}, fonttitle=\bfseries, sharp corners]
\begin{lstlisting}
while true {
    try {
        say "inside try";
        break;
    } finally {
        say "finally runs!";
    }
}
# Output:
# inside try
# finally runs!

# continue also respects finally
for i in 0..3 {
    try {
        if i == 1 { continue; }
        say "processing ", i;
    } finally {
        say "cleanup for ", i;
    }
}
# Output:
# processing 0 / cleanup for 0
# cleanup for 1
# processing 2 / cleanup for 2
\end{lstlisting}
\end{tcolorbox}

\subsection{Return Statement}

\texttt{return} exits a function with an optional value. It uses a sentinel error internally to unwind the call stack:

\begin{tcolorbox}[colback=codebg, colframe=darkpurple, colbacktitle=darkpurple, title=\textbf{\textcolor{codegold}{Return Statement}}, fonttitle=\bfseries, sharp corners]
\begin{lstlisting}
# Return with value
fn add(a, b) {
    return a + b;
}

# Early return from condition
fn abs(x) {
    if x >= 0 { return x; }
    return -x;
}

# Return without value (returns null)
fn do_nothing() {
    return;
}
\end{lstlisting}
\end{tcolorbox}

\begin{tcolorbox}[colback=lightlavender, colframe=softvioletdark, title=\textbf{\textcolor{darkpurple}{Design Note: Control Flow Signals}}, fonttitle=\bfseries, sharp corners]
\texttt{return}, \texttt{break}, and \texttt{continue} are implemented as sentinel errors (\texttt{RuntimeError::ReturnSignal}, \texttt{BreakSignal}, \texttt{ContinueSignal}) that unwind through the call stack. This ensures they propagate correctly through nested structures like \texttt{if}, \texttt{match}, \texttt{try/catch}, and \texttt{for} loops. Using \texttt{break} or \texttt{continue} outside a loop produces a runtime error.
\end{tcolorbox}
